\chapter{Retícula de firmas y carácter exponencial de masa}
\label{chap:masas}

La ley interna se formula primero sobre la retícula completa de firmas
\(\Lambda=\mathbb Z^5\). El
\cref{chap:extension-ruta-caracter} construye el extractor desde una extensión
superviviente, por su ruta observable y su registro, hasta \(\Lambda\); el
carácter formal precede a la especialización en los invariantes globales del
\cref{chap:torres}. El sector nulo queda fuera del codominio positivo.

La restricción histórica calibrada
\[
 \mathcal V_{10}=\{v_X:X\in\mathcal Q_{10}\}\subset\ZZ^5
\]
del \cref{app:tablas} constituye una aplicación particular de esa arquitectura,
no su dominio ni su catálogo. La asignación prospectiva de una identidad
física a una extensión superviviente permanece como interfaz separada. El
análisis exhaustivo de una reducción unifilar de 108 pasos determina con
precisión qué información pierde esa reducción y por qué no puede sustituir
al estado enriquecido.

\section{Dominio y factorización del carácter de masa}

Denotamos por
\[
 \Gamma_{108}^{\rm uni}
 :=I_9^2\times\{\pm1\}^2
\]
la familia finita de condiciones iniciales de la dinámica unifilar explícita
que se define en la sección siguiente. Este conjunto no se identifica con la
imagen del espacio completo de trayectorias TPK: establecer una proyección
sobreyectiva desde aquel espacio exigiría definir previamente el morfismo que
retiene exactamente estas cuatro coordenadas. Deben distinguirse las
aplicaciones
\[
 F_{108}^{\rm uni}:\Gamma_{108}^{\rm uni}\longrightarrow\ZZ^5
 \qquad
 L_{\rm tick}:\Gamma_{108}^{\rm enr}\longrightarrow\ZZ^5
 \qquad\text{y}\qquad
 \widehat\sigma_{10}:\mathcal Q_{10}\dashrightarrow\ZZ^5,
 \qquad X\longmapsto v_X,
\]
para las diez filas históricas consideradas. El primer mapa pertenece a la dinámica
unifilar finita declarada; el segundo es el extractor de registro enriquecido
definido y factorizado relativamente al cociente de registro en el
\cref{chap:extension-ruta-caracter}. El tercer mapa proporciona, con estatuto
de reconstrucción calibrada, las firmas enteras que coordinan
\(\mathcal V_{10}\). El carácter y el transductor actúan sobre estas últimas:
\[
 \mathcal Q_{10}
 \overset{\widehat\sigma_{10}}{\dashrightarrow}\ZZ^5
 \xrightarrow{\ \chi_{\rm mass}\ }\RR_{>0}
 \xrightarrow{\ \varsigma_{\rm MeV}\ }\mathcal U_E .
\]

\begin{figure}[htbp]
\centering
\begin{tikzpicture}[>=Latex,node distance=15mm and 20mm,
box/.style={draw=hmtblue,rounded corners,fill=hmtlight,align=center,
minimum height=11mm,text width=34mm}]
\node[box] (sp) {especie\\[-1mm]$X$};
\node[box,right=of sp] (v) {firma calibrada\\[-1mm]$v_X\in\ZZ^5$};
\node[box,below=of v] (m) {razón positiva\\[-1mm]$e^{\langle\Theta,v_X\rangle}$};
\node[box,left=of m] (u) {magnitud de energía\\[-1mm]$\varsigma_{\rm MeV}$};
\draw[->,densely dashed,thick] (sp)--node[above]{$\widehat\sigma_{10}$}(v);
\draw[->,thick] (v)--node[right]{$\chi_{\rm mass}$}(m);
\draw[->,thick] (m)--node[above]{$\varsigma_{\rm MeV}$}(u);
\end{tikzpicture}
\caption{Composición tipada de la firma entera, el carácter exponencial y el
transductor de energía. La primera flecha representa una reconstrucción
calibrada y se dibuja discontinua; la reducción unifilar se estudia como una
aplicación anterior con dominio propio.}
\label{fig:cadena-masa}
\end{figure}

\section{Extracción de firmas en la subdinámica unifilar de 108 pasos}

La familia \(\Gamma_{108}^{\rm uni}\) se parametriza por
\((r,c,h_0,v_0)\in I_9^2\times\{\pm1\}^2\) y produce una palabra de 108
dígitos. Es una dinámica finita autónoma; no se presupone aquí que agote una
imagen del espacio enriquecido \(\mathcal X_{\rm TPK}\). En el paso
\(t\),
\[
 \rho(t)=1+((t-1)\bmod9).
\]
Los residuos \(1,4,7\) activan el observable aditivo; los residuos \(2,5,8\),
el multiplicativo; \(3,6,9\) pertenecen al sector nulo. En los pasos
\(27,54,81\) se invierten las orientaciones \((h,v)\).

La recurrencia que determina la palabra es la siguiente. Sea
\(\operatorname{dr}_9(n)=1+((n-1)\bmod 9)\), con valores en
\(I_9=\{1,\ldots,9\}\), y escribamos
\((r_t,c_t,h_t,v_t)\) para el estado que precede a la lectura \(t\). A partir
de \((r_1,c_1,h_1,v_1)=(r,c,h_0,v_0)\), se aplican sucesivamente las reglas
\[
\begin{array}{c|c|c}
 \rho(t)&d_t&(r_t',c_t')\\ \hline
 1,4,7&\operatorname{dr}_9(r_t+c_t)&
       (r_t,\operatorname{dr}_9(c_t+h_t))\\
 2,5,8&\operatorname{dr}_9(r_tc_t)&
       (\operatorname{dr}_9(r_t+v_t),c_t)\\
 3,6,9&9&(r_t,c_t).
\end{array}
\]
Después de efectuar la lectura y el desplazamiento se define
\[
 (h_{t+1},v_{t+1})=
 \begin{cases}
   (-h_t,-v_t),&t\in\{27,54,81\},\\
   (h_t,v_t),&\text{en otro caso},
 \end{cases}
 \qquad
 (r_{t+1},c_{t+1})=(r_t',c_t').
\]
Estas ecuaciones determinan de manera única
\((d_1,\ldots,d_{108})\) para cada condición inicial de
\(\Gamma_{108}^{\rm uni}\).

Para la palabra producida, se define
\[
 n_A=\#\{t:d_t\in\{1,3,5,7\}\},
\]
\[
 s_C=\sum_{t:d_t\in\{2,4,6,8\}}\sigma_t,
 \qquad
 \sigma_t=(-1)^{\lfloor(t-1)/27\rfloor}.
\]
Al reagrupar en doce bloques de nueve pasos, numerados por
\(m\in\{0,\ldots,11\}\), definimos primero la función de signo dodecafásica
\[
 \Sigma_{12}(m)=(-1)^{\lfloor m/3\rfloor}
 =(1,1,1,-1,-1,-1,1,1,1,-1,-1,-1)_m.
\]
Sean entonces
\[
 s_C^{(m)}=\Sigma_{12}(m)\,
 \#\{\text{dígitos pares en el bloque }m\},
 \qquad
 \tau_m=\operatorname{sgn}s_C^{(m)}.
\]
Las restantes coordenadas son
\[
 k=\sum_{m:\tau_m=0}\Sigma_{12}(m),
\]
\[
 \nu_{120}=\sum_{C\in\{(0,4,8),(1,5,9),(2,6,10),(3,7,11)\}}
 \operatorname{sgn}\sum_{m\in C}s_C^{(m)},
\]
\[
 Q_3(\tau)=\sum_{m=0}^{11}\tau_m\tau_{m+3\bmod12},
 \qquad \nu_{270}^{Q_3}=Q_3(\tau).
\]

\begin{proposition}[Aplicación trayectoria--firma]
Las fórmulas anteriores definen una aplicación determinista
\[
 F_{108}^{\rm uni}:\Gamma_{108}^{\rm uni}\longrightarrow\ZZ^5,
 \qquad
 \gamma\longmapsto(n_A,s_C,k,\nu_{120},\nu_{270}^{Q_3}).
\]
\end{proposition}
\begin{proof}
Cada coordenada es un conteo o una suma finita de la palabra asociada a
\(\gamma\). Por consiguiente, el vector depende exclusivamente de la
trayectoria reducida y queda definido antes de introducir una masa o una
etiqueta de especie.
\end{proof}

\begin{theorem}[Obstrucción de imagen para la reducción unifilar]
La enumeración de \(\Gamma_{108}^{\rm uni}\) produce 53 vectores distintos y
\[
 \operatorname{im}F_{108}^{\rm uni}
 \cap\operatorname{im}\widehat\sigma_{10}=\varnothing .
\]
En particular, no existe una aplicación
\(\sigma_{\rm uni}:\mathcal Q_{10}
\to\Gamma_{108}^{\rm uni}\) que satisfaga
\[
 F_{108}^{\rm uni}\sigma_{\rm uni}
 =\widehat\sigma_{10}
\]
con la dinámica unifilar de esta sección. Esta obstrucción no excluye un
selector definido directamente sobre el estado TPK enriquecido.
\end{theorem}
\begin{proof}
El dominio parametrizado contiene
\(9\cdot9\cdot2\cdot2=324\) condiciones iniciales. Su imagen completa posee
53 vectores. Las proyecciones coordenadas son
\[
\begin{aligned}
n_A&\in\{8,14,20,24,28,30,34,44,50\},\\
s_C&\in\{-6,-4,-2,0,2,4,6\},\\
k&\in\{-4,-2,0,2,4\},\\
\nu_{120}&\in\{-4,-2,0,2,4\},\\
\nu_{270}^{Q_3}&\in\{-12,-8,-4,0\}.
\end{aligned}
\]
La comparación exhaustiva de los 53 vectores con los diez elementos de
\(\operatorname{im}\widehat\sigma_{10}\) da cero coincidencias. Si existiera
la factorización indicada, cada uno de esos diez elementos pertenecería a
\(\operatorname{im}F_{108}^{\rm uni}\), en contradicción con la intersección
vacía.
\end{proof}

El teorema localiza una obstrucción de la dinámica unifilar explícita. El
extractor sobre la trayectoria TPK enriquecida posee un dominio distinto,
pues conserva orientación, frontera, doce eventos y cociclos a nivel de paso;
su construcción se da en el \cref{chap:extension-ruta-caracter}. En la tabla
de contraste, las diez firmas ingresan mediante
\(\widehat\sigma_{10}\). La flecha física que debe fijarse de forma
prospectiva es
\[
 \mathcal Q_{10}\dashrightarrow\mathcal E_{108}^{\rm surv},
\]
no una búsqueda inversa de una ruta después de conocer la masa. Una vez
construida esa flecha, el extractor descendido
\(\overline L:\mathcal E_{108}^{\rm surv}\to\ZZ^5\) ya está tipado y
produce el morfismo de procedencia completo.

\section{Especialización global y carácter exponencial positivo}

Los coeficientes globales construidos en el \cref{chap:torres} fijan
\[
 \Theta:=\left(A_{\rm rad},\frac{\Cstarrad}{6},
 \Dfour,s_{120},\zCorr\right),
 \qquad \zCorr=135\Dfour^2.
\]
Para \(v\in\ZZ^5\), la evaluación del carácter formal en ese punto define
\[
 \chi_{\rm mass}(v)
 :=\operatorname{ev}_{\Theta}\bigl(\chi_{\rm form}(v;X)\bigr)
 =e^{\langle\Theta,v\rangle}>0.
\]

\begin{corollary}[Especialización del carácter de masa]
\label{cor:especializacion-caracter-masa}
\(\chi_{\rm mass}:(\ZZ^5,+)\to(\RR_{>0},\cdot)\) es un homomorfismo. Además,
\[
\chi_{\rm route}^{\rm mass}
:=\chi_{\rm mass}\circ L_{\rm tick}
\]
hace conmutar
\[
 \begin{array}{ccc}
 \Gamma_{108}^{\rm enr}&\xrightarrow{\ L_{\rm tick}\ }&\ZZ^5\\[1mm]
 \big\downarrow{\scriptstyle\chi_{\rm route}^{\rm mass}}&&
 \big\downarrow{\scriptstyle\chi_{\rm mass}}\\[1mm]
 \RR_{>0}&=&\RR_{>0}.
 \end{array}
\]
La salida sigue siendo adimensional y el resultado es exacto una vez fijado
\(\Theta\).
\end{corollary}
\begin{proof}
Es la composición del homomorfismo del
\cref{thm:caracter-formal-ruta} con la evaluación multiplicativa
\(\operatorname{ev}_{\Theta}\). En particular, para \(v,w\in\ZZ^5\),
\[
 \chi_{\rm mass}(v+w)=\chi_{\rm mass}(v)\chi_{\rm mass}(w).
\]
\end{proof}

Condicionada a una firma
\(v_X=(n_A,s_C,k,\nu_{120},\nu_{270})\), la salida adimensional construida es
\[
 \boxed{r_X^{(3)}:=\chi_{\rm mass}(v_X).}
\]
Una realización metrológica posterior define
\(E_X^{(3)}:=E_e\,r_X^{(3)}\) sobre la recta de energía. La igualdad con una
energía equivalente física no forma parte de la definición del carácter;
constituye el contraste de esa realización con una referencia de especie.

\section{Descomposición sucesiva de la acción de masa}

Para una firma fija, se introducen
\[
 E_0=n_AA_{\rm rad}+\frac{s_C}{6}\Cstarrad,
\]
\[
 E_1=E_0+k\Dfour,
 \qquad E_2=E_1+\nu_{120}s_{120},
 \qquad E_3=E_2+\nu_{270}\zCorr,
\]
\[
 \boxed{E_4=E_3+\nu_{180}s_{180}.}
\]
Con la escala electrónica interna, la restricción histórica calibrada
presenta los errores RMS relativos
\[
 2.2242372039\times10^{-3},\quad
 2.7620461148\times10^{-5},\quad
 4.7772213690\times10^{-6},
\]
\[
 4.7953121368\times10^{-7},\qquad
 \boxed{2.2292793330\times10^{-8}}.
\]

\begin{theorem}[Cuantización reconstructiva del residuo en la etapa \(E_4\)]
Para los diez elementos de la restricción \(\mathcal V_{10}\),
\[
 \nu_{180}(X)=\nint\!\left(
 \frac{\log(E_X^{\rm ref}/E_X^{(3)})}{s_{180}}
 \right),
\]
y, en consecuencia,
\[
 \left|\log\frac{E_X^{\rm ref}}{E_X^{(3)}}
 -\nu_{180}(X)s_{180}\right|\leq\frac{s_{180}}2.
\]
\end{theorem}
\begin{proof}
La primera igualdad se obtiene evaluando el cociente residual de cada firma.
La segunda es la desigualdad característica del entero más cercano.
\end{proof}

El vector
\[
 (2,7,-7,-1,-4,5,-4,-5,6,9)
\]
en el orden
\((\mu,\pi^0,\pi^+,K^+,K^0,p,n,\tau,W,Z)\) se obtuvo a partir del residuo
posterior a \(E_3\). La etapa \(E_4\) demuestra, por tanto, la cuantización de
ese residuo dentro de la malla \(s_{180}\); no constituye una selección
prospectiva de \(\nu_{180}\).

\section{Normalizaciones electrónicas y referencias de comparación}

La evaluación de calibración emplea la energía equivalente
\(E_e^{\rm ext}=0.510998950000000\,{\rm MeV}\); la evaluación interna usa
\(E_e^{\rm HMT}=0.510998922488066072\ldots\,{\rm MeV}\). Cada columna de la
tabla mantiene una sola normalización en todas sus etapas. Las referencias
metrológicas posteriores se incorporan como columnas de comparación y no
sustituyen retroactivamente los datos a partir de los cuales se obtuvieron las
firmas.
